% Captions of the figures, one \caption-ready macro per figure, each at most two short sentences:
% what is shown, and the one point to take from it; the decoding a reader needs (line styles,
% markers, scales) is in the sentence of the text that cites the figure. Notation: theory/CONVENTIONS.md.
% Colours are named by tone only ("dark", "light") and only where needed to read the figure, so the
% captions hold in greyscale print. Needs \usepackage{amssymb} (\mathfrak). The \ref labels are those
% of paper/jga/manuscript.tex, which reads this file and attaches \figcapOne ... \figcapEight to
% F1-F8 (\sref is its macro for a label of the electronic supplement). F9 belongs to the companion
% note, paper/arith/note.tex, which carries the same caption in its own notation. Every statement in
% a caption is asserted by the figure's script in figures/src/ or proved in the paper. Old and new
% word counts: paper/jga/FIGURE-RESTORE.md.

% F1 (Introduction) -- figures/out/F1.pdf
\newcommand{\figcapOne}{One triangle of $\mathcal O(2,8,8)$ (a) and of $\mathcal O(3,3,12)$ (b), in schematic shapes, coloured by $4\pi t\,\mathfrak h_t(x,x)$ at $t=0.02$. The two orbifolds share their first two heat invariants, yet heat lingers differently near their cone points.}

% F2 (Section 5) -- figures/out/F2.pdf
\newcommand{\figcapTwo}{Exact tilings of the hyperboloid by the triangle groups of $\mathcal O(2,8,8)$ (a) and $\mathcal O(3,3,12)$ (b), in an oblique orthographic view. The coloured triangle and its pale mirror image form one copy of each orbifold.}

% F3 (Section 3) -- figures/out/F3.pdf
\newcommand{\figcapThree}{The mirror argument of Theorem~\ref{thm:sigsep} for $(2,8,8)$, $(3,3,12)$ (a) and $(1;15)$, $(0;3,3,5,5)$ (b). Discs mark $X^*$ and rings $-X^*$; they coincide only for an orbifold against itself.}

% F4 (Section 3) -- figures/out/F4.pdf
\newcommand{\figcapFour}{The largest $K_{\mathrm{mult}}$ on complete equal-area classes and on constructed pairs, against $s=\mathrm{Area}/2\pi$. The open growth question lies between the two step curves.}

% F5 (Section 7) -- figures/out/F5.pdf
\newcommand{\figcapFive}{(a) The difference of the heat traces of $\mathcal O(2,8,8)$ and $\mathcal O(3,3,12)$ from computed spectra, against its cone-point term. (b) Differences of heat traces in the family $(0;3,3,3,3)$, against the predicted geodesic term.}

% F6 (Section 4, what heat does not hear) -- figures/out/F6.pdf
\newcommand{\figcapSix}{(a) Four members of the family $(0;3,3,3,3)$ with their shortest closed geodesics; (b) its eigenvalues below $120$ at eight moduli. All members share every heat invariant, yet the spectrum moves.}

% F7 (Section 5) -- figures/out/F7.pdf
\newcommand{\figcapSeven}{Reciprocal sums $R$ of the hyperbolic triads of sum $S$, one band for each least order $p=2,\dots,8$. Adjacent bands first overlap at the dots and first collide at the circles.}

% F8 (Section 6) -- figures/out/F8.pdf
\newcommand{\figcapEight}{Recovery error of the cone orders against a worst-case error $\delta$ in every heat invariant, for $(2,3,7)$, $(2,8,8)$ and $(4,4,4)$. The slopes are the exponents of Theorem~\ref{thm:S3}.}

% F9 (companion note, Section 6) -- figures/out/F9.pdf; the note's own copy uses \N and \ciso.
\newcommand{\figcapNine}{(a) The real locus of $C_{27/2}$ in the chart $x+y+z=1$, with the permutations of
$(1,4,4)$ and $(1,1,4)$, that is, $(2,8,8)$ and $(3,3,12)$ up to scale. (b) The number $\mathcal N(S)$ of pairs of sum at most $S$ sharing $c_1,c_2$, divided by $S$
(heavy), and its fit (dashed), far above the isosceles lower bound, which grows like
$(c_{\mathrm{iso}}+o(1))\log S$ (thin).}

% E1-E3: figures of the eigen manuscript (paper/eigen/manuscript.tex), same conventions.
% E1 (Section 8) -- figures/out/E1.pdf
\newcommand{\figcapEOne}{(a) The eigenvalues $\lambda_1,\dots,\lambda_6$ of $\mathcal O(2,3,m)$ against $m$, with their upper bounds; (b) the same eigenvalues, rescaled as $h_m^2(\lambda_j-\frac14)/\pi^2$. As the cone point of order $m$ becomes cusp-like, the low spectrum crowds towards $\frac14$.}

% E2 (Section 6) -- figures/out/E2.pdf
\newcommand{\figcapETwo}{The gaps between the signature terms of the heat trace of $\mathcal O(2,8,8)$ and of its five competitors, with the error of the geodesic term and the error of keeping finitely many eigenvalues. Where half the nearest gap exceeds both errors, finitely many eigenvalues decide the signature.}

% E3 (Section 7) -- figures/out/E3.pdf
\newcommand{\figcapEThree}{Eigenvalues needed to decide the signature on computed spectra, observed and certified, against the a-priori count of Theorem~\ref{thm:E}, by systole. The certificate needs $10^2$ to $10^{11}$ times fewer eigenvalues.}
